\scanpage{55}
\uline{Polygône combinatoire et repères}

\textcircled{\small 1.} Polygône combinatoire $(S,A,R)$,
\diagram{\begin{tikzcd}[ampersand replacement=\&,row sep=tiny,column sep=0pt]
{}\&\text{sommets}\&{}\&\text{arêtes}\&{}\&\text{repères}\\
S\arrow[ur]{},\&{}\& A\arrow[ur]{}\,\text{ens,}\&{}\& R\arrow[ur]{}\subset S\times A
\end{tikzcd}}
relation, axiomes \quad 1) $R\to A$ de degré 2 \quad 2) $R\to S$ de
degré 2, \quad 3) Connexité. [Quand on laisse tomber
l'axiome 3) de connexité, on parle de ``contour
combinatoire'' ; quand on laisse tomber \deleted{\unsure{$R'$|$R$|$R\subset$}} l'axiome 2)
($R\to S$ de degré 2) on parle de graphe ; quand
on \deleted{prend} \inserted{se donne} $R\to S$ et $R\to A$, sans supposer $R\hookrightarrow S\times A$,
on parle de \deleted{\unsure{graphe|graphes}} complexe \deleted{combinatoire} \inserted{1-cellulaire}
combinatoire (ou parfois encore de ``graphe'' en
un sens légèrement généralisé).

\textcircled{\small 2.} Si
\diagram{\begin{tikzcd}[ampersand replacement=\&,row sep=0pt,column sep=large]
\raisebox{-0.9ex}{$R$}\arrow[r]\& S\\
\phantom{R}\arrow[r]\& A
\end{tikzcd}}
est un contour, on définit
\[
\begin{aligned}
\sigma_1&\quad(\text{involution sans pt fixe de }R\text{ tel que }R/\sigma_1\xrightarrow{\sim}S)\\
\sigma_0&\quad(\qquad-\qquad-\qquad\qquad-\qquad-\quad R/\sigma_0\simeq A)
\end{aligned}
\]
Soit $\unsure{G_1|\mathbb G_1|\mathfrak G_1}$ le groupe défini par
\[
G_1=\{\sigma_0,\sigma_1\mid\sigma_0^2=\sigma_1^2=1\}\xrightarrow{\sim}\mathbb Z/2\mathbb Z*\mathbb Z/2\mathbb Z
\]
alors on a un foncteur allant de la
catégorie des contours dans la cat. des $G_1$-ens.
$R$ tels que \quad a) $\sigma_0,\sigma_1$ opèrent ss pt fixe dans $E$
b) $\forall x\in E$, on a $\sigma_0x\ne\sigma_1x$ [ou encore les $G_1$-orbites
dans $E$ sont de cardinal $\unsure{>|\geq|\to}2$] -- qui exprime
que $R\to S\times A$ (où $S\overset{\mathrm{dfn}}=R/\sigma_1$, $A\overset{\mathrm{dfn}}=R/\sigma_0$) est
% archive sheet number: 54
